% !PRINTED 395
\weberpage{395}
Wir bilden das reelle syzygetische Complextripel, in dem alle
Doppeltangenten vorkommen müssen:
\[
\begin{array}{rl}
1) & x\,y,\ p\,q \\
2) & x\,p,\ y\,q \\
3) & x\,q,\ y\,p.
\end{array}
\]
Ist eine fünfte reelle Doppeltangente $r$ vorhanden, so können wir
annehmen, sie komme im Complexe 1) vor. Dann muss aber dieser Complex
ein drittes reelles Paar $r\,s$ enthalten, und entweder ein viertes
reelles Paar, in welchem Falle der Satz 3. schon zutrifft, oder ein
conjugirtes Paar $\varrho\varrho'$.

Enthält der Complex 2) nicht lauter reelle Doppeltangenten, ein Fall,
in dem gleichfalls der Satz 3. zutreffen würde, so muss in 2) ein
imaginäres Paar $\xi\eta$ vorkommen, und danach betrachten wir also die
folgenden Complexe:
\[
\begin{array}{rl}
1) & x\,y,\ p\,q,\ r\,s,\ \varrho\varrho' \\
2) & x\,p,\ y\,q,\ \xi\eta \\
4) & x\,r,\ y\,s.
\end{array}
\]
Die beiden Complexe 2), 4) sind aber \gesperrt{azygetisch}, da z. B. $r$ in 2)
nicht vorkommt [weil es in 1) vorkommt]. Folglich kann 4) auch nur
eine der beiden $\xi, \eta$ enthalten, und folglich können, da der Complex
4) reell ist, $\xi, \eta$ nicht conjugirt sein. Es ergiebt sich daraus für
2) und 4) je ein conjugirtes Doppelpaar, und wir haben:
\[
\begin{array}{rl}
2) & x\,p,\ y\,q,\ \xi\eta,\ \xi'\eta' \\
4) & x\,r,\ y\,s,\ \xi\zeta,\ \xi'\zeta'.
\end{array}
\]
worin $\zeta, \zeta'$ wieder ein Paar conjugirter Doppeltangenten bedeutet.

Nun bilden wir den sowohl mit 2) als mit 4) syzygetischen reellen
Complex
\[
5)\quad \xi\xi',\ \eta\eta',\ \zeta\zeta',
\]
der keines der Elemente $x, y, p, q, r, s$ enthalten kann, und ferner die
beiden conjugirt imaginären Complexe
\[
\begin{array}{rl}
6) & x\xi,\ p\eta,\ r\zeta \\
7) & x\xi',\ p\eta',\ r\zeta',
\end{array}
\]
und 5), 6), 7) bilden ein \gesperrt{azygetisches Tripel}. $y, q, s$ kommen in 6)
und 7) nicht vor, denn sonst müssten sie mit je zweien der $x, p, r$
azygetisch sein, was nach 2) und 4) unmöglich ist.
